\section{Introduction}
\label{sec:intro}

A proactive streaming assistant watches a video as it arrives and must decide, at
every moment, whether to speak. Almost universally that decision is implemented as
a scalar confidence compared against a threshold, and almost as universally the
threshold is tuned per task, because different tasks plainly want different
eagerness: an alert that fires once per video is not the same object as a running
commentary. Tuning nine thresholds for nine tasks and reporting the resulting gain
is a standard move.

This paper performs that fit carefully for Qwen2.5-Omni-7B on OmniPro Online, and
reports that it does not work --- along with the measurements that establish
\emph{why} it does not, and the one intervention on the same gate that does.

We make three contributions.

\paragraph{An audit that the fitted quantity is the used quantity.} Before fitting,
we found that the shipped gate ignored the parameter we intended to fit: its fire
test was a fixed global $0.5$, and the per-task threshold table, the gate mode and
the refractory period were inert. We also found that $p_{\mathrm{hit}}$ is high on
the first tick of every video --- the model answers the gating question from the
prompt before it has seen any of the video --- which under an edge-triggered gate
with a long refractory consumed an entire video's emission budget at $t=1.0$\,s.
Both are repaired before any number in this paper is measured. We report them
because a fit of an inert parameter is the failure mode that a well-run grid search
is least able to detect on its own.

\paragraph{A test of whether a fitted threshold is a measurement.} A grid of ten
cells scored on fifteen videos always has a maximum. We test each selection by a
paired bootstrap over videos that replays the actual selection rule, and report a
re-selection rate against chance alongside a confidence interval on the difference
from the best rival. Under this test no task's fitted threshold beats its own best
alternative outside noise, and the fits re-select themselves 25--61\,\% of the time
--- one of them at 6\,\%, below its 8\,\% chance rate. Two apparent fits dissolve
entirely on inspection: one task scores exactly $0.000$ at all ten thresholds, and
another's entire threshold curve is one matched event, or zero, out of sixteen.

\paragraph{A mechanism, and a positive result.} The reason is measurable rather
than conjectural: $p_{\mathrm{hit}}$ does not rank event-adjacent ticks above quiet
ones. Per-task AUC is $0.431$--$0.551$ with every 95\,\% interval containing $0.5$,
and this survives a $\pm 10$\,s sweep of the tick clock, a widened $\pm 10$\,s
tolerance, and a check that the score is not degenerate (232--284 distinct values
spanning $0.000$--$0.997$). A threshold applied to a score that cannot say
\emph{when} can only buy emission \emph{volume} --- and volume has exactly one
useful setting. Pooled over 135 videos the single global operating point is
identified and stable, re-selected in 94\,\% of resamples, and moving it from the
shipped gate's effective $0.5$ is the one change here with a reproducible effect.
Split nine ways over fifteen videos each, the same quantity is not identifiable at
all.

The result is therefore not that thresholds do not matter. It is that this gate has
one identifiable degree of freedom rather than nine, and that the difference
between those two statements is invisible to any protocol that reports an argmax
without asking whether it recurs.
